% !TeX root = booklet-03.tex
\chapter*{مرور تعادل جریان و ابزار تحلیل مشخصه انتقال}
\addcontentsline{toc}{chapter}{مرور تعادل جریان و ابزار تحلیل مشخصه انتقال}
\section*{از صفر و یک تا \lr{Inverter}}
در یک مدار منطقی، صفر و یک دو مقدار قراردادی برای نمایش اطلاعات‌اند.
مدار الکتریکی، این مقدارها را با ولتاژ نمایش می‌دهد. در منطق مثبت،
ولتاژ نزدیک زمین معمولاً صفر و ولتاژ نزدیک تغذیه مثبت معمولاً یک
تفسیر می‌شود. زمین، مرجع \(\qty{0}{\volt}\) است و تغذیه را
\(V_{DD}>0\) می‌نامیم. مرز دقیق بازه‌های معتبر، موضوع
\lr{Logic Levels} و \lr{Noise Margin} در ادامه بخش \secnum{۲.۵}
است \cite[مقدمه بخش \secnum{۲.۵} و اسلاید ۳]{weste,dcslides}.

\lr{Inverter} تابع \lr{NOT} را پیاده می‌کند. خروجی منطقی آن
متمم ورودی است. جدول \ref{tab:inverter-truth} این تابع را نشان
می‌دهد. در یک \lr{CMOS Inverter}، ترانزیستور \lr{pMOS} مسیر
بالاکش به تغذیه و ترانزیستور \lr{nMOS} مسیر پایین‌کش به زمین را
فراهم می‌کند \cite[شکل \secnum{۲.۲۵}]{weste}.
\begin{table}[!htbp]
\centering
\caption{جدول درستی \lr{Inverter} در منطق مثبت. مقدارها منطقی‌اند.}
\label{tab:inverter-truth}
\begin{tabular}{cc}
\toprule
ورودی & خروجی\\
\midrule
۰ & ۱\\
۱ & ۰\\
\bottomrule
\end{tabular}
\end{table}

وقتی ورودی پایین است، \lr{nMOS} خاموش و \lr{pMOS} روشن است.
خروجی از راه \lr{pMOS} به تغذیه نزدیک می‌شود. وقتی ورودی بالا
است، \lr{pMOS} خاموش و \lr{nMOS} روشن است و خروجی به زمین نزدیک
می‌شود. در ورودی‌های میانی، هر دو ترانزیستور ممکن است روشن باشند.
برای تعیین ولتاژ میانی، باید مقدار جریان‌ها را محاسبه کنیم.

\section*{ولتاژ پایانه، جریان و تعادل گره}
نماد \(V_{AB}=V_A-V_B\) اختلاف ولتاژ پایانه اول نسبت به پایانه
دوم است. مثلاً اگر \(V_G=\qty{0.5}{\volt}\) و
\(V_S=\qty{1.2}{\volt}\) باشند، اختلاف ولتاژ چنین محاسبه می‌شود.
\[
 V_{GS}=0.5-1.2=-\qty{0.7}{\volt}.
\]
با جابه‌جایی دو
پایانه، علامت عوض می‌شود و \(V_{SG}=\qty{0.7}{\volt}\) خواهد بود.
این نکته در \lr{pMOS} برای جلوگیری از اشتباه علامت اهمیت دارد.

جریان، نرخ عبور بار و واحد آن آمپر است.
\lr{Kirchhoff's Current Law (KCL)} بیان می‌کند که جمع جریان‌های
ورودی به یک گره با مجموع جریان‌های خروجی، با احتساب تغییر بار
ذخیره‌شده، سازگار است. در خروجی \lr{Inverter}، بزرگی جریان
بالاکش را \(I_p\ge0\)، جریان پایین‌کش را \(I_n\ge0\) و جریان
بار خروجی را \(I_L\) می‌نامیم. قرارداد \(I_L>0\) یعنی بار از
خروجی جریان می‌گیرد. اگر ظرفیت ثابت خروجی \(C_L\) باشد، داریم
\[
 C_L\frac{\mathrm{d}V_{\mathrm{out}}}{\mathrm{d}t}
 =I_p-I_n-I_L.
\]
حاصل‌ضرب ظرفیت با واحد فاراد در نرخ تغییر ولتاژ با واحد ولت بر
ثانیه، واحد آمپر دارد. اگر \(I_p>I_n+I_L\)، بار خروجی افزایش
می‌یابد و ولتاژ بالا می‌رود. اگر طرف راست منفی باشد، ولتاژ کاهش
می‌یابد. در حالت ماندگار، مشتق ولتاژ صفر است. برای بار بدون جریان
\lr{DC}، \(I_L=0\) و شرط ماندگار به \(I_p=I_n\) تبدیل می‌شود
\cite[صفحه ۸۸]{weste}\cite[اسلاید ۳]{dcslides}.

ورودی گیت بعدی می‌تواند بار عمدتاً خازنی ایجاد کند. چنین باری در
مدل ایده‌آل، هنگام ثابت‌بودن ولتاژ جریان پیوسته نمی‌گیرد، اما زمان
رسیدن خروجی به مقدار نهایی را تغییر می‌دهد. یک مقاومت متصل به
خروجی می‌تواند جریان \lr{DC} بگیرد و باید در \(I_L\) وارد شود.
پس فرض بار در تحلیل باید مشخص باشد.

\section*{تابع، نمودار و شیب منفی}
رابطه \(y=f(x)\) یعنی برای ورودی \(x\)، مقدار خروجی \(y\) از یک
قاعده محاسبه می‌شود. در مشخصه انتقال، محور افقی
\(V_{\mathrm{in}}\) و محور عمودی \(V_{\mathrm{out}}\) است.
یک نقطه \((\qty{0.5}{\volt},\qty{1.1}{\volt})\) یعنی در ورودی
\(\qty{0.5}{\volt}\)، خروجی \(\qty{1.1}{\volt}\) است.
این نقطه صرفاً برای خواندن محورهای نمودار فرض شده است.

شیب متوسط بین دو نقطه، نسبت تغییر خروجی به تغییر ورودی است.
مثلاً اگر ورودی \(\qty{0.01}{\volt}\) زیاد و خروجی
\(\qty{0.20}{\volt}\) کم شود، شیب متوسط برابر
\((-0.20)/(0.01)=-20\) است. علامت منفی نشان می‌دهد افزایش
ورودی با کاهش خروجی همراه است. هرچه بزرگی شیب بیشتر باشد، تغییر
کوچک ورودی، تغییر بزرگ‌تری در خروجی می‌دهد.
مشتق \(\mathrm{d}V_{\mathrm{out}}/\mathrm{d}V_{\mathrm{in}}\)
شیب محلی منحنی در یک نقطه است. واحد این نسبت ولت بر ولت و در
نتیجه بی‌بعد است. این شیب در ادامه با \lr{Gain} ارتباط پیدا می‌کند.

\section*{مربع، ریشه و حل معادله درجه دوم}
در رابطه جریان مدل \lr{Shockley}، مربع ولتاژ ظاهر می‌شود. برای
حل، سه اتحاد زیر کاربرد دارند
\[
 (x-y)^2=x^2-2xy+y^2,\qquad
 (x-y)(x+y)=x^2-y^2,\qquad
 (x-y)^2\ge0.
\]
اگر \(z^2=d\) و \(d>0\)، دو جواب جبری
\(z=\pm\sqrt d\) وجود دارند. نماد \(\sqrt d\) به‌تنهایی ریشه
نامنفی را نشان می‌دهد. در معادله \(z^2-2bz+c=0\)، با افزودن
\(b^2\) به شکل مناسب می‌توان مربع کامل ساخت
\[
 z^2-2bz=-c
 \quad\Longrightarrow\quad
 (z-b)^2=b^2-c
 \quad\Longrightarrow\quad
 z=b\pm\sqrt{b^2-c}.
\]
برای جواب حقیقی، باید \(b^2-c\ge0\) باشد. در مدار، شرط‌های ناحیه
کاری و بازه ولتاژ نیز باید برقرار باشند. همین شرط‌ها تعیین می‌کنند
کدام ریشه جبری قابل استفاده است.

برای نمونه، \(z^2-0.8z+0.04=0\) دو ریشه تقریبی
\(0.05359\) و \(0.74641\) دارد. اگر مسئله شرط \(0\le z\le0.4\)
داشته باشد، فقط ریشه نخست در دامنه قرار می‌گیرد. در مثال‌های مدار،
کمیت \(z\) می‌تواند ولتاژ باشد و در آن صورت ضریب‌ها هم واحدهای
متناسب خواهند داشت.

